\subsection{\texorpdfstring{第~\ref{sec:levi-civita-connection}~节}{Levi-Civita 联络}}\label{subsec:solutions-levi-civita-connection}

\begin{description}[font=\normalfont,leftmargin=2em,labelwidth=1.5em]
\RGsolution{ex:isometry-preserves-levi-civita}{isometry-preserves-levi-civita}
\begin{enumerate}[label=(\alph*),ref=\alph*,leftmargin=*]
\item\label{sol:isometry-preserves-connection-by-koszul} 利用 $f$ 是等距映射及 Koszul 公式, 有
\begin{align*}
  2\langle f_*\nabla_XY,f_*Z\rangle_{f(p)}
  &=2\langle\nabla_XY,Z\rangle_p=X_p\cdot\langle Y,Z\rangle+Y_p\cdot\langle Z,X\rangle\\
  &\quad-Z_p\cdot\langle X,Y\rangle-\langle[Y,Z],X\rangle_p
        -\langle[X,Z],Y\rangle_p-\langle[Y,X],Z\rangle_p.
\end{align*}
另一方面,
\begin{align*}
  &2\langle\widetilde\nabla_{f_*X}f_*Y,f_*Z\rangle_{f(p)}\\
  &\quad=(f_*X)_{f(p)}\cdot\langle f_*Y,f_*Z\rangle
        +(f_*Y)_{f(p)}\cdot\langle f_*Z,f_*X\rangle\\
  &\qquad-(f_*Z)_{f(p)}\cdot\langle f_*X,f_*Y\rangle-\langle[f_*Y,f_*Z],f_*X\rangle_{f(p)}\\
  &\qquad-\langle[f_*X,f_*Z],f_*Y\rangle_{f(p)}-\langle[f_*Y,f_*X],f_*Z\rangle_{f(p)}\\
  &\quad=X_p\cdot(\langle f_*Y,f_*Z\rangle\circ f)+Y_p\cdot(\langle f_*Z,f_*X\rangle\circ f)\\
  &\qquad-Z_p\cdot(\langle f_*X,f_*Y\rangle\circ f)-\langle f_*[Y,Z],f_*X\rangle_{f(p)}\\
  &\qquad-\langle f_*[X,Z],f_*Y\rangle_{f(p)}-\langle f_*[Y,X],f_*Z\rangle_{f(p)}\\
  &\quad=X_p\cdot\langle Y,Z\rangle+Y_p\cdot\langle Z,X\rangle
    -Z_p\cdot\langle X,Y\rangle-\langle[Y,Z],X\rangle_p\\
  &\qquad-\langle[X,Z],Y\rangle_p-\langle[Y,X],Z\rangle_p
    =2\langle f_*\nabla_XY,f_*Z\rangle_{f(p)}.
\end{align*}
\item 设 $c:I\to M$ 是测地线, 考虑映射 $\widetilde c:=f\circ c:I\to N$. 由于
\[
  \dot{\widetilde c}(t)=\frac{\mathrm{d}(f\circ c)}{\mathrm{d}t}
   =(\mathrm{d}f)_{c(t)}\dot c,
\]
有
\[
  \widetilde\nabla_{\dot{\widetilde c}}\dot{\widetilde c}
   =\widetilde\nabla_{(\mathrm{d}f)_{c(t)}\dot c}\bigl((\mathrm{d}f)_{c(t)}\dot c\bigr)
   \overset{(\ref{sol:isometry-preserves-connection-by-koszul})}{=}
     (\mathrm{d}f)_{c(t)}(\nabla_{\dot c}\dot c)=0,
\]
因此 $\widetilde c$ 是 $N$ 中的测地线.
\end{enumerate}

\RGsolution{ex:spherical-geodesics-and-transport}{spherical-geodesics-and-transport}
\begin{enumerate}[label=(\alph*),ref=\alph*,leftmargin=*]
\item 因为 $\frac{\partial}{\partial\theta}$ 是固定 $\varphi$ 而改变 $\theta$ 得到的 $S^2$ 上曲线的切向量, 所以
\[
  \frac{\partial}{\partial\theta}=\frac{\partial\phi}{\partial\theta}
   =(\cos\theta\cos\varphi,\cos\theta\sin\varphi,-\sin\theta),
\]
因而
\[
  g_{\theta\theta}=\left\langle\frac{\partial}{\partial\theta},\frac{\partial}{\partial\theta}\right\rangle=1.
\]
类似地,
\[
  \frac{\partial}{\partial\varphi}=\frac{\partial\phi}{\partial\varphi}
   =(-\sin\theta\sin\varphi,\sin\theta\cos\varphi,0),
\]
故
\begin{align*}
  g_{\varphi\varphi}
  &=\left\langle\frac{\partial}{\partial\varphi},\frac{\partial}{\partial\varphi}\right\rangle=\sin^2\theta;\\
  g_{\theta\varphi}=g_{\varphi\theta}
  &=\left\langle\frac{\partial}{\partial\theta},\frac{\partial}{\partial\varphi}\right\rangle=0.
\end{align*}
因此, $\mathbb{R}^3$ 的 Euclid 度量在 $S^2$ 上诱导的度量为
\[
  g=\mathrm{d}\theta\otimes\mathrm{d}\theta+\sin^2\theta\,\mathrm{d}\varphi\otimes\mathrm{d}\varphi.
\]
\item 有
\[
  (g_{ij})=\begin{pmatrix}g_{\theta\theta}&g_{\theta\varphi}\\g_{\varphi\theta}&g_{\varphi\varphi}\end{pmatrix}
   =\begin{pmatrix}1&0\\0&\sin^2\theta\end{pmatrix},
\]
因而
\[
  (g^{ij})=(g_{ij})^{-1}=\begin{pmatrix}1&0\\0&\dfrac{1}{\sin^2\theta}\end{pmatrix}.
\]
由这些矩阵可以容易地算出 Christoffel 符号. 例如,
\begin{align*}
  \Gamma^\theta_{\varphi\varphi}
  &=\frac12\sum_{i=1}^2g^{\theta i}
    \left(\frac{\partial g_{\varphi i}}{\partial\varphi}
         +\frac{\partial g_{\varphi i}}{\partial\varphi}
         -\frac{\partial g_{\varphi\varphi}}{\partial x^i}\right)\\
  &=\frac12 g^{\theta\theta}\left(0+0-\frac{\partial(\sin^2\theta)}{\partial\theta}\right)
    =-\sin\theta\cos\theta.
\end{align*}
八个 Christoffel 符号中只有三个非零: 上面算出的一个及
\[
  \Gamma^\varphi_{\theta\varphi}=\Gamma^\varphi_{\varphi\theta}=\cot\theta.
\]
\item 测地线方程为
\begin{align*}
  \ddot\theta+\sum_{i,j=1}^2\Gamma^\theta_{ij}\dot{x}^i\dot{x}^j=0
  &\iff \ddot\theta-\sin\theta\cos\theta\,\dot\varphi^2=0;\\
  \ddot\varphi+\sum_{i,j=1}^2\Gamma^\varphi_{ij}\dot{x}^i\dot{x}^j=0
  &\iff\ddot\varphi+2\cot\theta\,\dot\theta\dot\varphi=0.
\end{align*}
坐标表达式为 $\widehat c(t)=(\theta(t),\varphi(t))=(\frac{\pi}{2},t)$ 的曲线 $c$ 显然是这些方程的一个解. 因此赤道 $\theta=\frac{\pi}{2}$ 是一条测地线的像.
\item 绕 $\mathbb{R}^3$ 中任一过原点的轴的旋转, 都是保持 $S^2$ 的 $\mathbb{R}^3$ 等距映射. 由于这里考虑的 $S^2$ 上的度量由 $\mathbb{R}^3$ 上的 Euclid 度量诱导, 可见这样的旋转确定 $S^2$ 上的一个等距映射.
\item 给定点 $p\in S^2$ 和向量 $v\in T_pS^2$, 存在旋转 $R:\mathbb{R}^3\to\mathbb{R}^3$, 使得 $R(p)=(1,0,0)$ 且 $R(v)=(0,1,0)$. 具有这些初始条件的测地线显然是坐标表达式为 $\widehat c(t)=(\theta(t),\varphi(t))=(\frac{\pi}{2},t)$ 的曲线 $c$, 其像为赤道. 由 Exercise~\ref{ex:levi-civita-connection}(\ref{ex:isometry-preserves-levi-civita}), 以 $v\in T_pS^2$ 为初始条件的测地线必为 $R^{-1}\circ c$. 由于 $c$ 的像是 $S^2$ 与平面 $z=0$ 的交集, $R^{-1}\circ c$ 的像就是 $S^2$ 与某个过原点的平面的交集, 即一个大圆.
\item 例如顶点为 $(1,0,0)$, $(0,1,0)$ 和 $(0,0,1)$ 的三角形 (参见 Fig.~\ref{fig:three-right-angles-on-sphere}).
\RGfigure[.42\textwidth]{three-right-angles-on-sphere}{$S^2$ 上有三个直角的测地三角形}
\item\label{sol:parallel-transport-along-a-latitude} 平行移动的方程为
\begin{align*}
  \dot V^\theta+\sum_{i,j=1}^2\Gamma^\theta_{ij}\dot{x}^iV^j=0
  &\iff\dot V^\theta+\Gamma^\theta_{\varphi\varphi}V^\varphi=0\\
  &\iff\dot V^\theta-\sin\theta_0\cos\theta_0V^\varphi=0
\end{align*}
(因为沿 $c$ 有 $\dot\varphi=1$) 和
\begin{align*}
  \dot V^\varphi+\sum_{i,j=1}^2\Gamma^\varphi_{ij}\dot{x}^iV^j=0
  &\iff\dot V^\varphi+\Gamma^\varphi_{\varphi\theta}V^\theta=0\\
  &\iff\dot V^\varphi+\cot\theta_0V^\theta=0
\end{align*}
(因为沿 $c$ 有 $\dot\varphi=1$). 这些方程蕴含
\[
  \ddot V^\theta+\cos^2\theta_0V^\theta=0
  \iff V^\theta=A\cos((\cos\theta_0)t)+B\sin((\cos\theta_0)t),
\]
其中 $A,B\in\mathbb{R}$ 为常数, 因而
\[
  V^\varphi=\frac{1}{\sin\theta_0\cos\theta_0}\dot V^\theta
  =-\frac{A}{\sin\theta_0}\sin((\cos\theta_0)t)+\frac{B}{\sin\theta_0}\cos((\cos\theta_0)t).
\]
初始条件为 $V^\theta(0)=1$, $V^\varphi(0)=0$, 故 $A=1$, $B=0$, 从而
\begin{align*}
  V^\theta&=\cos((\cos\theta_0)t);\\
  V^\varphi&=-\frac{1}{\sin\theta_0}\sin((\cos\theta_0)t).
\end{align*}
特别注意,
\[
  \langle V(t),V(t)\rangle=(V^\theta)^2+\sin^2\theta_0(V^\varphi)^2=1.
\]
所以 $V(0)$ 和 $V(2\pi)$ 之间的角 $\alpha$ 由
\[
  \cos\alpha=\langle V(0),V(2\pi)\rangle=V^\theta(2\pi)=\cos(2\pi\cos\theta_0)
\]
给出, 即
\[
  \alpha=2\pi\cos\theta_0\quad\text{或}\quad\alpha=2\pi(1-\cos\theta_0)
\]
(取决于选择度量哪一个角).
\item 利用 $S^2$ 上任一点都可以由适当的等距映射送到 $(0,0,1)$, 只需证明 $(1,0,0)$ 在 $S^2$ 中没有开邻域 $U\subset S^2$ 等距于带 Euclid 度量的开集 $V\subset\mathbb{R}^2$. 注意任一这样的邻域都包含 (\ref{sol:parallel-transport-along-a-latitude}) 中所给曲线 $c(t)$ 的像 (取充分小的 $\theta_0>0$). 若 $U$ 与 $W$ 等距, 则 $U$ 上的 Levi-Civita 联络将是平凡联络, 从而 (\ref{sol:parallel-transport-along-a-latitude}) 中的平行向量场 $V(t)$ 将满足 $V(0)=V(2\pi)$. 因为对任意 $\theta_0\in(0,\frac{\pi}{2})$ 这都不成立, 所以 $U$ 与 $W$ 不可能等距.
\item 平行公设在 $S^2$ 中不成立, 因为任意两条测地线的像或者重合, 或者恰好在两点相交.
\end{enumerate}

\RGsolution{ex:hyperbolic-plane}{hyperbolic-plane}
\begin{enumerate}[label=(\alph*),ref=\alph*,leftmargin=*]
\item 如 Exercise~\ref{ex:riemannian-manifolds}(\ref{ex:left-invariant-metric}) 所示, 对所有 $(x,y)\in H$ 和 $v,w\in T_{(x,y)}H$, 必有
\[
  \langle v,w\rangle_{(x,y)}=
  \left\langle(\mathrm{d}L_{(x,y)^{-1}})_{(x,y)}v,(\mathrm{d}L_{(x,y)^{-1}})_{(x,y)}w\right\rangle_e.
\]
由第~\ref{chap:differentiable-manifolds}~章的 Exercise~\ref{ex:lie-groups}(\ref{ex:affine-maps-lie-group}), 知
\[
  (\mathrm{d}L_{(x,y)})_{(z,w)}=\begin{pmatrix}y&0\\0&y\end{pmatrix}
\]
和
\[
  (x,y)^{-1}=\left(-\frac{x}{y},\frac{1}{y}\right).
\]
因此有
\begin{align*}
  &\left\langle v^x\frac{\partial}{\partial x}+v^y\frac{\partial}{\partial y},
    w^x\frac{\partial}{\partial x}+w^y\frac{\partial}{\partial y}\right\rangle_{(x,y)}\\
  &\quad=\left\langle\frac1y v^x\frac{\partial}{\partial x}+\frac1y v^y\frac{\partial}{\partial y},
        \frac1y w^x\frac{\partial}{\partial x}+\frac1y w^y\frac{\partial}{\partial y}\right\rangle_{(0,1)}\\
  &\quad=\frac1{y^2}\left\langle v^x\frac{\partial}{\partial x}+v^y\frac{\partial}{\partial y},
    w^x\frac{\partial}{\partial x}+w^y\frac{\partial}{\partial y}\right\rangle_{(0,1)}
   =\frac1{y^2}(v^xw^x+v^yw^y),
\end{align*}
即
\[
  g=\frac1{y^2}(\mathrm{d}x\otimes\mathrm{d}x+\mathrm{d}y\otimes\mathrm{d}y).
\]
\item 有
\[
  (g_{ij})=\begin{pmatrix}g_{xx}&g_{xy}\\g_{yx}&g_{yy}\end{pmatrix}
  =\begin{pmatrix}\frac1{y^2}&0\\0&\frac1{y^2}\end{pmatrix},
\]
因此
\[
  (g^{ij})=(g_{ij})^{-1}=\begin{pmatrix}y^2&0\\0&y^2\end{pmatrix}.
\]
由这些矩阵可以容易地算出 Christoffel 符号. 例如,
\begin{align*}
  \Gamma^x_{xy}
  &=\frac12\sum_{i=1}^2g^{xi}
    \left(\frac{\partial g_{yi}}{\partial x}+\frac{\partial g_{xi}}{\partial y}
      -\frac{\partial g_{xy}}{\partial x^i}\right)\\
  &=\frac12g^{xx}\left(0+\frac{\partial}{\partial y}\left(\frac{1}{y^2}\right)-0\right)
  =-\frac1y.
\end{align*}
八个 Christoffel 符号中只有四个非零: 上面算出的一个及
\[
  \Gamma^x_{yx}=-\Gamma^y_{xx}=\Gamma^y_{yy}=-\frac1y.
\]
\item 测地线方程为
\begin{align*}
  \ddot x+\sum_{i,j=1}^2\Gamma^x_{ij}\dot{x}^i\dot{x}^j=0
  &\iff\ddot x-\frac2y\dot x\dot y=0;\\
  \ddot y+\sum_{i,j=1}^2\Gamma^y_{ij}\dot{x}^i\dot{x}^j=0
  &\iff\ddot y+\frac1y\dot x^2-\frac1y\dot y^2=0.
\end{align*}
直接代入可检验曲线 $\alpha$ 和 $\beta$ 满足这些方程. 还显然有 $\alpha(\mathbb{R})$ 是 $y$ 轴的正半轴. 由于
\[
  \tanh^2t+\frac1{\cosh^2t}=1,
\]
所以 $\beta(\mathbb{R})$ 是单位圆与 $H$ 的交集.
\item 由 Exercise~\ref{ex:levi-civita-connection}(\ref{ex:isometry-preserves-levi-civita}), 等距映射把测地线的像送到测地线的像. 由于 $H$ 上的度量左不变, 任意左平移都是等距映射. 特别地, 对所有 $a\in\mathbb{R}$ 和 $b>0$,
\[
  L_{(a,1)}(x,y)=(x+a,y)
\]
与
\[
  L_{(0,b)}(x,y)=(bx,by)
\]
都是等距映射. 因而所有竖直半直线和所有圆心在 $x$ 轴上的半圆都是 $H$ 中测地线的像 (参见 Fig.~\ref{fig:hyperbolic-plane-semicircles}). 另一方面, 给定 $p\in H$ 和 $v\in T_pH$, $v$ 总与其中之一相切. 实际上, 若 $v$ 是竖直的, 则它与经过 $p$ 的竖直半直线相切. 若 $v$ 不是竖直的, 则它与以 $x$ 轴和过 $p$ 且垂直于 $v$ 的直线的交点为圆心的半圆相切 (参见 Fig.~\ref{fig:hyperbolic-plane-semicircles}). 所以, 以 $v\in T_pM$ 为初始条件的测地线的像或者是竖直半直线, 或者是圆心在 $x$ 轴上的半圆.
\RGfigure[.77\textwidth]{hyperbolic-plane-semicircles}{(a) 双曲平面的测地线及平行公设不成立. (b) 与给定向量相切的测地线. (c) 过两点的测地线. (d) 测地四边形的内角}
\item 若 $p$ 和 $q$ 位于同一条竖直半直线上, 则不存在包含 $p$ 和 $q$ 且圆心在 $x$ 轴上的半圆. 若 $p$ 和 $q$ 不在同一条竖直半直线上, 则存在唯一一个包含 $p$ 和 $q$ 且圆心在 $x$ 轴上的半圆, 圆心是以 $p$ 和 $q$ 为端点的线段的垂直平分线与 $x$ 轴的交点 (参见 Fig.~\ref{fig:hyperbolic-plane-semicircles}).
\item 在球面 $S^2$ 上, 经过任意一对对径点的测地线有无穷多条, 且彼此不是重新参数化所得. 在带通常 Euclid 度量的 $\mathbb{R}^2\setminus\{0\}$ 中, 不存在连接点 $(1,0)$ 和 $(-1,0)$ 的测地线.
\item 因为双曲平面的度量是 Euclid 度量乘以一个函数, 用两种度量测得的角相等. 考虑由两条竖直半直线和两个不同半圆形成的测地四边形, 两个半圆的圆心都是其中一条半直线与 $x$ 轴的交点 (参见 Fig.~\ref{fig:hyperbolic-plane-semicircles}). 容易检验此四边形的内角和小于 $2\pi$. 而每个开集 $U\subset H$ 都包含一个这样的四边形. 若 $U$ 等距于 $V\subset\mathbb{R}^2$, 则此四边形的内角和必恰为 $2\pi$. 因此 $U$ 不可能与 $V$ 等距.
\item 平行公设在双曲平面中不成立. 相反, 容易看出, 给定测地线 $c:\mathbb{R}\to H$ 和点 $p\notin c(\mathbb{R})$, 存在无穷多条测地线 $\widetilde c:\mathbb{R}\to H$ (不计重新参数化), 使得 $p\in\widetilde c(\mathbb{R})$ 且 $c(\mathbb{R})\cap\widetilde c(\mathbb{R})=\varnothing$ (参见 Fig.~\ref{fig:hyperbolic-plane-semicircles}).
\end{enumerate}

\RGsolution{ex:induced-connection-and-projection}{induced-connection-and-projection}
\begin{enumerate}[label=(\alph*),ref=\alph*,leftmargin=*]
\item 首先注意, 若 $p\in N$, 则 $(\widetilde\nabla_{\widetilde X}\widetilde Y)_p$ 只依赖于 $\widetilde X_p=X_p$ 以及 $\widetilde Y$ 在一条与 $X_p$ 相切的曲线上的值, 因而可把这条曲线选在 $N$ 上. 由于在 $N$ 上 $\widetilde Y=Y$, 可知 $\widetilde\nabla_{\widetilde X}\widetilde Y|_N$ 只依赖于 $X$ 和 $Y$. 定义 $\nabla:\mathfrak{X}(N)\times\mathfrak{X}(N)\to\mathfrak{X}(N)$ 为
\[
  \nabla_XY=(\widetilde\nabla_{\widetilde X}\widetilde Y)^\top.
\]
下面证明 $\nabla$ 是 $N$ 上的 Levi-Civita 联络. $\nabla$ 在 $N$ 上定义仿射联络, 这是因为 $\widetilde\nabla$ 是 $M$ 上的仿射联络, 且正交投影 $\vphantom{w}^\top:TM|_N\to TN$ 是线性的. 另一方面, 给定 $X,Y,Z\in\mathfrak{X}(N)$, 有
\begin{align*}
  X\cdot\langle Y,Z\rangle
  &=\widetilde X\cdot\langle\widetilde Y,\widetilde Z\rangle
    =\langle\widetilde\nabla_{\widetilde X}\widetilde Y,\widetilde Z\rangle
     +\langle\widetilde Y,\widetilde\nabla_{\widetilde X}\widetilde Z\rangle\\
  &=\langle(\widetilde\nabla_{\widetilde X}\widetilde Y)^\top,\widetilde Z\rangle
     +\langle\widetilde Y,(\widetilde\nabla_{\widetilde X}\widetilde Z)^\top\rangle\\
  &=\langle\nabla_XY,Z\rangle+\langle Y,\nabla_XZ\rangle,
\end{align*}
其中使用了如下事实: 若 $v\in T_pN$, 则对所有 $w\in T_pM$, 有 $\langle v,w\rangle=\langle v,w^\top\rangle$. 所以 $\nabla$ 与诱导度量相容. 最后,
\begin{align*}
  \nabla_XY-\nabla_YX
  &=(\widetilde\nabla_{\widetilde X}\widetilde Y)^\top-
      (\widetilde\nabla_{\widetilde Y}\widetilde X)^\top\\
  &=(\widetilde\nabla_{\widetilde X}\widetilde Y-\widetilde\nabla_{\widetilde Y}\widetilde X)^\top
    =[\widetilde X,\widetilde Y]^\top=[X,Y],
\end{align*}
这里使用了 $N$ 上 $[\widetilde X,\widetilde Y]=[X,Y]$ 的事实 [参见第~\ref{chap:differentiable-manifolds}~章的 Exercise~\ref{ex:vector-fields}(\ref{ex:lie-bracket-tangent-to-submanifold})]. 所以 $\nabla$ 是对称的, 因而是 $\langle\langle\cdot,\cdot\rangle\rangle$ 的 Levi-Civita 联络.
\item 若 $(M,\langle\cdot,\cdot\rangle)$ 是带 Euclid 度量的 $\mathbb{R}^3$, $N\subset\mathbb{R}^3$ 是曲面, $c:I\to N$ 是曲线, 则有
\[
  \nabla_{\dot c}\dot c=(\widetilde\nabla_{\dot c}\dot c)^\top=\ddot c^\top.
\]
因此, $c$ 是 $N$ 的测地线, 当且仅当其加速度与 $N$ 正交. 假设 $c$ 同时是平面 $L$ 上的一条曲线, 且 $L$ 与 $N$ 正交 (即对每个 $p\in L\cap N$, $T_pL$ 与 $T_pN$ 正交). 那么 $\dot c(t)$ 和 $\ddot c(t)$ 都在 $T_{c(t)}L$ 中. 因为 $T_{c(t)}L\cap T_{c(t)}N=\mathbb{R}\dot c(t)$, 在这种情形下, $\ddot c^\top$ 是 $\dot c$ 的倍数. 最后, 若以弧长为 $c$ 的参数, 则
\[
  \langle\dot c,\dot c\rangle=1\implies\langle\ddot c,\dot c\rangle=0\implies\ddot c^\top=0.
\]
因此, 任何以弧长为参数且像落在 $L\cap N$ 上的曲线 $c$ 都是 $N$ 的测地线. 类似地, 不难证明, 若 $L$ 不与 $N$ 正交且 $\ddot c$ 不为零, 则 $c$ 不是 $N$ 的测地线. 由此, 下列结论显然成立.
\begin{enumerate}[label=(\roman*),ref=\roman*,leftmargin=*]
\item 在球面 $S^2$ 上, 任一大圆 (即 $S^2$ 与过原点的平面的任一交集) 都是测地线的像. 任一不是大圆的圆 (即 $S^2$ 与不过原点的平面的任一交集) 都不是测地线的像.
\item 旋转环面与任何对称平面 (即关于该平面的反射保持环面不变的平面) 的交集都是测地线的像 [另见 Exercise~\ref{ex:minimizing-properties-of-geodesics}(\ref{ex:fixed-point-submanifold-of-isometry})]. 旋转环面与任何不是对称平面的平面的交集都不是测地线的像.
\item 圆锥的母线都是测地线的像 (也可由它们本身已是 $\mathbb{R}^3$ 的测地线看出). 用垂直于圆锥轴的平面截圆锥所得的圆不是测地线的像.
\item 任意旋转曲面与任何包含旋转轴的平面的交集都是测地线的像. 注意这包含上面所有情形.
\end{enumerate}
\item 这直接来自
\[
  \nabla_{\dot c}V=(\widetilde\nabla_{\dot c}V)^\top=\dot V^\top
\]
只依赖于算子 $\vphantom{w}^\top$ 这一事实. 若两个曲面沿曲线 $c$ 相切, 则此算子对两个曲面相同.
\item 圆锥的一个参数化 $\varphi:(0,+\infty)\times(0,2\pi)\to\mathbb{R}^3$ 由
\[
  \varphi(r,\theta)=(r\cos\theta,r\sin\theta,r\cot\alpha)
\]
给出, 其中 $\alpha$ 为母线与轴的夹角. 此时可以求得诱导度量为
\[
  g=\frac{1}{\sin^2\alpha}\,\mathrm{d}r\otimes\mathrm{d}r+r^2\,\mathrm{d}\theta\otimes\mathrm{d}\theta.
\]
定义新坐标
\[
  \begin{cases}
    r':=\dfrac r{\sin\alpha},\\[2pt]
    \theta':=\theta\sin\alpha,
  \end{cases}
\]
有
\[
  g=\mathrm{d}r'\otimes\mathrm{d}r'+r'^2\,\mathrm{d}\theta'\otimes\mathrm{d}\theta',
\]
这说明去掉一条母线的圆锥与 $\mathbb{R}^2$ 中由 $\theta'\in(0,2\pi\sin\alpha)$ 给出的开集等距. 特别地, 如图所示, 绕圆锥平行移动一周, 初始向量与最终向量的夹角为 $2\pi(1-\sin\alpha)$. 若球面上的圆由
\[
  c(t)=(\sin\theta_0\cos t,\sin\theta_0\sin t,\cos\theta_0)
\]
参数化, 则由初等几何知, 沿 $c$ 与球面相切的圆锥满足 $\alpha=\frac{\pi}{2}-\theta_0$ (参见 Fig.~\ref{fig:cone-tangent-to-sphere}). 因此初始向量与最终向量的夹角为 $2\pi(1-\cos\theta_0)$ [参见 Exercise~\ref{ex:levi-civita-connection}(\ref{ex:spherical-geodesics-and-transport})].
\RGfigure[.33\textwidth]{cone-tangent-to-sphere}{$\alpha$ 与 $\theta_0$ 的关系}
\end{enumerate}
\end{description}
